\documentclass[11pt]{amsart}
\usepackage[margin=1in]{geometry}
\usepackage[T1]{fontenc}
\usepackage[utf8]{inputenc}
\usepackage{amsmath,amssymb,amsthm,amscd}
\usepackage{comment}
\usepackage[colorlinks=true,linkcolor=blue,citecolor=blue,urlcolor=blue]{hyperref}
\setlength{\parindent}{0pt}
\setlength{\parskip}{0.4\baselineskip}
\newcommand{\Z}{\mathbb Z}
\newcommand{\Q}{\mathbb Q}
\newcommand{\R}{\mathbb R}
\newcommand{\C}{\mathbb C}
\newcommand{\T}{\mathbb T}
\newcommand{\coker}{\operatorname{coker}}
\newcommand{\rank}{\operatorname{rank}}
\newcommand{\diag}{\operatorname{diag}}
\newcommand{\Res}{\operatorname{Res}}
\newcommand{\Rep}{R_{\C}}
\newcommand{\degc}{\operatorname{deg}}
\newcommand{\KC}[2]{K^{#1}_{#2}(\underline E#2)}
\newcommand{\KH}[2]{K_{#1}^{#2}(\underline E#2)}
\newtheorem{theorem}{Theorem}[section]
\newtheorem{proposition}[theorem]{Proposition}
\newtheorem{lemma}[theorem]{Lemma}
\theoremstyle{remark}
\newtheorem{remark}[theorem]{Remark}
\numberwithin{equation}{section}
\title[Torsion in crystallographic group $C^*$-algebras]{Counterexamples to torsion-freeness in the
$K$-theory of crystallographic group $C^*$-algebras}
\author{Cristian Sarmiento}
\address{Departamento de Matem\'aticas\\Universidad Nacional de Colombia, sede Bogot\'a\\Cra. 30 cll 45 - Ciudad Universitaria\\Bogot\'a, Colombia}
\email{cdsarmientos@unal.edu.co}
\author{Mario Vel\'asquez}
\address{Departamento de Matem\'aticas\\Universidad Nacional de Colombia, sede Bogot\'a\\Cra. 30 cll 45 - Ciudad Universitaria\\Bogot\'a, Colombia}
\email{mavelasquezme@unal.edu.co}
\date{\today}

\begin{document}
\begin{abstract}
We give a negative answer to a question of S\'anchez Salda\~na--Vel\'asquez and L\"uck asking whether the topological K-theory of the reduced group C*-algebra of a split crystallographic group with finite cyclic
holonomy is always torsion-free. For two explicit split
crystallographic groups with translation subgroup of rank four and holonomy of order four, we compute the equivariant K-theory  of the classifying space for proper actions and find nontrivial two-torsion. The equivariant universal coefficient theorem transfers this torsion to proper equivariant K-homology, and the Baum--Connes conjecture identifies the resulting groups with the topological K-theory of the corresponding reduced group C*-algebras.
\end{abstract}
\subjclass[2020]{Primary 46L80; secondary 20H15, 19K33}
\keywords{Crystallographic groups, equivariant $K$-theory, group $C^*$-algebras, Baum--Connes conjecture}
\maketitle

\section{Introduction}

The computation of the topological $K$-theory of group $C^*$-algebras of crystallographic groups has been studied by several authors under various hypotheses on the holonomy action. Davis and L\"uck \cite{DL} considered groups $\Z^n\rtimes C_p$, where $p$ is prime and the conjugation action of $C_p$ on $\Z^n$ is free outside the origin. They computed the real and complex topological $K$-theory of the corresponding group $C^*$-algebras by computing the equivariant $K$-homology of the classifying spaces for proper actions and using the Baum--Connes conjecture. Langer and L\"uck \cite{LL} extended this calculation from prime-order holonomy to arbitrary finite cyclic holonomy $C_m$, under the assumption that the action is free outside the origin.

More recently, S\'anchez Salda\~na and Vel\'asquez \cite{SV} removed the freeness hypothesis when the holonomy order $m$ is square-free. In particular, they proved that
\[
 K_*(C_r^*(\Z^n\rtimes C_m))
\]
is a finitely generated torsion-free abelian group for every action of $C_m$ on $\Z^n$ when $m$ is square-free. They asked whether the same torsion-freeness conclusion holds for arbitrary finite cyclic holonomy. L\"uck formulates the same problem as Question~17.50 on page~548 of \cite{LuckBook}. The examples below give a negative answer: for two actions of $C_4$ on $\Z^4$, the $K$-theory of the reduced group $C^*$-algebra contains nontrivial $2$-torsion. 

S\'anchez Salda\~na and Vel\'asquez observed that counterexamples to a cohomological conjecture of Adem--Ge--Pan--Petrosyan might provide a source of counterexamples to torsion-freeness. We follow precisely this idea. The two order-four integral representations used here were exhibited by Petrosyan and Putrycz \cite{PP} in their study of split crystallographic groups for which the collapsing of the canonical Lyndon--Hochschild--Serre spectral sequence fails.

Our calculation is done through two reductions. First, each four-dimensional group is written as $\Gamma_j=H\rtimes_{\alpha_j}\Z$, where $H$ is a common three-dimensional crystallographic group. Viewing the corresponding classifying space for proper actions as a mapping torus gives a Mayer--Vietoris sequence, which reduces the calculation for $\Gamma_j$ to the equivariant $K$-theory of $H$ and the action induced by $\alpha_j$. Second, the group $H$ is expressed as the amalgam
\[
 H=p4*_{p2}p4
\]
of two planar crystallographic groups over a common planar subgroup. The Mayer--Vietoris sequence for this amalgam reduces the computation of $K_H^*(\underline EH)$ to the restriction map from $p4$ to $p2$. We compute this restriction using explicit equivariant line bundles on the square two-torus. The equivariant universal coefficient theorem and Baum--Connes conjecture then identify the above computations with the $K$-theory groups of the reduced group $C^*$-algebras.

We now state the two examples and fix our conventions. Let

\[
 A_1=\begin{pmatrix}-1&0&0&0\\0&0&0&-1\\0&1&0&1\\0&0&-1&1\end{pmatrix},
 \qquad
 A_2=\begin{pmatrix}0&1&0&0\\-1&0&0&1\\0&0&-1&1\\0&0&0&1\end{pmatrix},
 \qquad \Gamma_j=\Z^4\rtimes_{A_j}C_4.
\]
These order-four matrices occur in Remark~3 and equation~(3.1), respectively, of \cite{PP}.

We write $K_G^i(Y)$ for complex equivariant topological \emph{$K$-cohomology} of a proper $G$-space, and $K_i^G(Y)$ for proper equivariant \emph{$K$-homology}. In degree zero, on a finite proper $G$-CW-complex, the former is the Grothendieck group of finite-dimensional equivariant complex vector bundles \cite{LO}.

\begin{theorem}\label{thm:main}
For the classifying space for proper actions $\underline E\Gamma_j=\R^4$,
\begin{equation}\label{eq:cohom-results}
\begin{array}{c|cc}
 &K^0_{\Gamma_j}(\R^4)&K^1_{\Gamma_j}(\R^4)\\ \hline
 \Gamma_1&\Z^8\oplus\Z/2&\Z^8\oplus\Z/2\\
 \Gamma_2&\Z^8\oplus\Z/2&\Z^8.
\end{array}
\end{equation}
Consequently,
\begin{equation}\label{eq:hom-results}
\begin{array}{c|cc}
 &K_0^{\Gamma_j}(\R^4)&K_1^{\Gamma_j}(\R^4)\\ \hline
 \Gamma_1&\Z^8\oplus\Z/2&\Z^8\oplus\Z/2\\
 \Gamma_2&\Z^8&\Z^8\oplus\Z/2.
\end{array}
\end{equation}
The second table also gives $K_i(C_r^*(\Gamma_j))$, since the
Baum--Connes assembly map is an isomorphism for these groups.
\end{theorem}

\section{Reduction to planar stabilizers}

Put
\[
 R=\begin{pmatrix}0&-1\\1&0\end{pmatrix},\qquad
 H=\Z^3\rtimes_{\diag(R,-1)}C_4.
\]
Let $u,v,z$ be the three canonical generators of $\Z^3$ and $s$ the generator of $C_4$ expressed as the order four rotation in the first two coordinates and the reflection for the last one. Thus
\[
 sus^{-1}=v,\qquad svs^{-1}=u^{-1},\qquad szs^{-1}=z^{-1}.
\]
In order to relate the groups $\Gamma_j$ with $H$ let us consider the integer changes of lattice basis
\[
 B_1=\begin{pmatrix}0&0&1&0\\1&0&0&1\\-1&1&0&0\\0&1&0&0\end{pmatrix},
 \qquad B_2=\diag(1,-1,1,1)
\]
give
\[
 B_1^{-1}A_1B_1=
 \begin{pmatrix}0&-1&0&-1\\1&0&0&0\\0&0&-1&0\\0&0&0&1\end{pmatrix},
 \qquad
 B_2^{-1}A_2B_2=
 \begin{pmatrix}0&-1&0&0\\1&0&0&-1\\0&0&-1&1\\0&0&0&1\end{pmatrix}.
\]

The last coordinate is fixed by the holonomy and defines a surjection $\Gamma_j\to\Z$ with kernel $H$. Conjugation by the fourth element of the canonical basis yields
\begin{equation}\label{eq:alphas}
 \Gamma_j=H\rtimes_{\alpha_j}\Z,\qquad
 \alpha_1(s)=us,\qquad \alpha_2(s)=vz^{-1}s,
\end{equation}
and both automorphisms fix $u,v,z$.

Write points of $\R^2\times\R$ as $(x,y)$, with $x\in\R^2$ and $y\in\R$. Project the affine action of $H$ onto the last coordinate. On this line, $u,v$ act trivially, $z$ acts by $y\mapsto y+1$, and $s$ by $y\mapsto-y$. With vertices at $\frac12\Z$, the quotient is an edge $[0,1/2]$. Its stabilizers are
\[
 V_0=\langle u,v,s\rangle\cong p4,\qquad
 V_1=\langle u,v,zs\rangle\cong p4,\qquad
 E=\langle u,v,s^2\rangle\cong p2,
\]
where
\[
 p4=\Z^2\rtimes_R C_4,\qquad
 p2=\Z^2\rtimes_{-I}C_2.
\]
There are no edge inversions. The associated graph of groups gives $H=V_0*_E V_1$. Since $(zs)^2=s^2$, both edge inclusions are the standard inclusion $p2\hookrightarrow p4$ under these identifications.

Replacing each vertex and edge by its classifying space of proper actions produces
$\underline EH=\R^2\times\R$, with
\[
 s(x,y)=(Rx,-y)
\]
and $u,v,z$ acting by translations. 

Recall that a proper $G$-$CW$-complex $Y$ is a model for $\underline EG$ if and only if $Y^F$ is contractible for every finite subgroup $F\leq G$ \cite[Theorem~1.9]{LuckSurvey}.
The affine actions considered here are proper and cocompact. A finite subgroup fixes the average of an orbit, and its fixed-point set is an affine subspace, hence contractible. Thus these are models for the classifying spaces for proper actions. They admit equivariant triangulations with finitely many cell orbits, as does the original $\Gamma_j$-action on $\R^4$.

\subsection{\texorpdfstring{$K$}{K}-theory of planar groups}
\label{sec:planar}

Let
\[
 X=\C/(\Z+i\Z),\qquad C_4=\langle s\rangle,
 \qquad C_2=\langle s^2\rangle.
\]
With $C_4$ acting on $X$ by $ s[w]=[iw]$, for every $w\in\C$. 
Taking the quotient by the freely acting normal subgroup $\Z^2$ we may identify
\begin{equation}\label{eq:descent}
 K_{p d}^i(\R^2)\cong K_{C_d}^i(X),\qquad d=2,4.
\end{equation}
 The same argument identifies $K_{\Gamma_j}^i(\R^4)$ with $K_{C_4}^i(\R^4/\Z^4)$ for the original action of $A_j$.

The quotient $X/C_2$ is a sphere with four points of isotropy order two, represented by
\[
 x_{00}=[0],\quad x_{10}=[1/2],\quad
 x_{01}=[i/2],\quad x_{11}=[(1+i)/2].
\]
The quotient $X/C_4$ is a sphere with isotropy orders $4,4,2$. Its exceptional orbits are
\[
 \mathsf p=x_{00},\qquad \mathsf q=x_{11},\qquad
 \{x_{10},x_{01}\}.
\]
%For example, these descriptions follow from the planar domains $[0,1]\times[0,1/2]$ for $p2$ and $[0,1/2]^2$ for $p4$, with the rotation identifications on their boundaries.

With the actions of $C_2,C_4$ defined above, let $q_d\colon X\to X/C_d\cong S^2$, for $d=2,4$. We equip $X/C_d$ with the unique orientation for which $q_d$ is locally orientation-preserving away from the exceptional orbits. 

Recall that for a complex virtual line bundle $\chi$ over an oriented closed surface $Y$, its Chern number is the integer
\[
 \langle c_1(\chi),[Y]\rangle,
\]
where $c_1(\chi)\in H^2(Y;\Z)$ is the first Chern class and $[Y]$ is the fundamental class determined by the orientation. Let $\mathcal H\to S^2$ be the dual Hopf line bundle, normalized by $\langle c_1(\mathcal H),[S^2]\rangle=1$, and put
\[
 \beta_d=q_d^*([\mathcal H]-1).
\]
The pullback $q_d^*\mathcal{H}$ has the canonical $C_d$-action
$g\cdot(x,v)=(gx,v)$ for $g\in C_d$, where
$v\in\mathcal{H}_{q_d(x)}$. This is well-defined because
$q_d(gx)=q_d(x)$; in particular, every stabilizer of $x$ acts trivially
on its fiber. We give the trivial line bundle in $[\mathcal{H}]-1$ the
trivial $C_d$-action as well, so $\beta_d$ defines a class in
$K^0_{C_d}(X)$. 

The degree will record the topological information that is not detected by
the representations on the exceptional fibers. The following
Lemma makes this precise for equivariant
$K$-theory classes and will be used to compute the restriction map from
$p4$ to $p2$.
\begin{lemma}\label{lem:surface-basis}
For these two actions, $K_{C_2}^1(X)=K_{C_4}^1(X)=0$, and fiber restriction gives exact sequences
\begin{equation}\label{eq:fiber-sequence}
 0\longrightarrow\Z\beta_d\longrightarrow K_{C_d}^0(X)
 \xrightarrow{\mathrm{fib}}\mathcal L_d\longrightarrow0,
 \qquad d=2,4,
\end{equation}
where
\[
\begin{aligned}
 \mathcal L_2&=\{(\sigma_{00},\sigma_{10},\sigma_{01},\sigma_{11})
 \in\Rep(C_2)^4:\dim\sigma_{ab}\text{ is constant}\},\\
 \mathcal L_4&=\{(\sigma_{\mathsf p},\sigma_{\mathsf q},\sigma_R)
 \in\Rep(C_4)^2\oplus\Rep(C_2):
 \dim\sigma_{\mathsf p}=\dim\sigma_{\mathsf q}=\dim\sigma_R\}.
\end{aligned}
\]
Moreover, $\degc\beta_2=2$ and $\degc\beta_4=4$. Consequently a class in either group is determined by its fiber representations at the exceptional orbits and its degree.
\end{lemma}

\begin{proof}
Choose pairwise disjoint closed disks around the four exceptional points of $X/C_2$, or the three exceptional orbits of $X/C_4$, and let $Y_d$ be the complement of their interiors. Choose slightly larger open disks and let $B_d$ be their union. Let $A_d$ be an open neighborhood of $Y_d$ such that $A_d\cap B_d$ is the union of one annular collar for each boundary component. Thus $S^2=A_d\cup B_d$, and $A_d$ deformation retracts onto $Y_d$.

Let $U_d=q_d^{-1}(A_d)$, $V_d=q_d^{-1}(B_d)$ and $W_d=U_d\cap V_d$. The space $U_d$ is free and each component of $V_d$ retracts equivariantly onto the corresponding exceptional orbit. The intersection $W_d$ is a union of free annuli. Since $U_d$ and $W_d$ are free $C_d$-spaces, passage to orbit spaces identifies their equivariant $K$-theory with ordinary $K$-theory. The Mayer--Vietoris sequences are therefore
\begin{align}
 K_{C_2}^0(X)&\longrightarrow
 \Z\oplus\Rep(C_2)^4
 \xrightarrow{r_0^{(2)}}\Z^4
 \longrightarrow K_{C_2}^1(X)
 \longrightarrow\Z^3
 \xrightarrow{r_1^{(2)}}\Z^4
 \longrightarrow K_{C_2}^0(X),
 \label{eq:MV-planar-C2}\\
 K_{C_4}^0(X)&\longrightarrow
 \Z\oplus\Rep(C_4)^2\oplus\Rep(C_2)
 \xrightarrow{r_0^{(4)}}\Z^3
 \longrightarrow K_{C_4}^1(X)
 \longrightarrow\Z^2
 \xrightarrow{r_1^{(4)}}\Z^3
 \longrightarrow K_{C_4}^0(X).
 \label{eq:MV-planar-C4}
\end{align}
Here the initial $\Z$ records the common rank of a bundle over $Y_d$. After ordering the boundary circles compatibly with the exceptional orbits, the degree-zero maps are
\begin{align*}
 r_0^{(2)}(n;\sigma_{00},\sigma_{10},\sigma_{01},\sigma_{11})
 &=\bigl(n-\dim\sigma_{00},n-\dim\sigma_{10},
         n-\dim\sigma_{01},n-\dim\sigma_{11}\bigr),\\
 r_0^{(4)}(n;\sigma_{\mathsf p},\sigma_{\mathsf q},\sigma_R)
 &=\bigl(n-\dim\sigma_{\mathsf p},
         n-\dim\sigma_{\mathsf q},n-\dim\sigma_R\bigr).
\end{align*}
Both maps are surjective, and their kernels are precisely $\mathcal L_2$ and $\mathcal L_4$. In suitable integral bases, restriction from $Y_d$ to its oriented boundary gives
\[
 r_1^{(2)}\colon\Z^3\longrightarrow\Z^4,\quad
 (a,b,c)\longmapsto(a,b,c,-a-b-c),
 \qquad
 r_1^{(4)}\colon\Z^2\longrightarrow\Z^3,\quad
 (a,b)\longmapsto(a,b,-a-b).
\]
Both maps are injective and have cokernel $\Z$. Exactness of \eqref{eq:MV-planar-C2}--\eqref{eq:MV-planar-C4} first gives $K_{C_2}^1(X)=K_{C_4}^1(X)=0$, and then gives
\[
 0\longrightarrow\coker r_1^{(d)}
 \longrightarrow K_{C_d}^0(X)
 \longrightarrow\ker r_0^{(d)}\longrightarrow0.
\]
Thus this is \eqref{eq:fiber-sequence}, once its left-hand generator is identified.

Let $\ell_2=4$ and $\ell_4=3$, and consider the summation map
\[
 \Sigma_d\colon\Z^{\ell_d}\longrightarrow\Z,
 \qquad (t_1,\ldots,t_{\ell_d})\longmapsto
 t_1+\cdots+t_{\ell_d}.
\]
The displayed formulas for $r_1^{(d)}$ show that $\operatorname{im}r_1^{(d)}=\ker\Sigma_d$. Hence $\Sigma_d$ induces an explicit isomorphism
\begin{equation*}
 \overline\Sigma_d\colon\coker r_1^{(d)}\xrightarrow{\cong}\Z.
\end{equation*}
In particular, the class of any standard basis vector of $\Z^{\ell_d}$ is a generator of the cokernel.

It remains to identify the image of this generator under the Mayer--Vietoris connecting homomorphism. Apply non-equivariant $K$-theory to the open cover $S^2=A_d\cup B_d$, formed by a neighborhood $A_d$ of $Y_d$ and the union $B_d$ of the slightly enlarged disks around the exceptional points. The degree-one
restriction-difference map is $r_1^{(d)}$. Hence the following Mayer--Vietoris connecting homomorphism factors through the cokernel, giving
\[
 \partial_{S^2}\colon\coker r_1^{(d)}\longrightarrow
 \widetilde K^0(S^2).
\]
With the boundary orientations chosen above, it sends the class of the first standard basis vector to $[\mathcal H]-1\in\widetilde K^0(S^2)$. Concretely, this class is obtained by using a clutching function of winding number one on the first boundary circle and the constant clutching function on the remaining circles; its first Chern number is one, so it is the class of the dual Hopf line bundle with our normalization.

Pullback by $q_d$ is compatible with the two Mayer--Vietoris sequences. On each free annulus, quotient identifies the equivariant boundary group upstairs with the ordinary boundary group downstairs. Writing $\partial_d$ for the equivariant connecting homomorphism, viewed as a map on the corresponding cokernel, the naturality square is
\[
\begin{CD}
 \coker r_1^{(d)} @>{\partial_{S^2}}>>
     \widetilde K^0(S^2)\\
 @V{\mathrm{id}}VV @VV{q_d^*}V\\
 \coker r_1^{(d)} @>{\partial_d}>> K_{C_d}^0(X).
\end{CD}
\]
Therefore
\[
 \partial_d([e_1])
 =q_d^*([\mathcal H]-1)=\beta_d.
\]
Thus the injection of $\coker r_1^{(d)}$ into $K_{C_d}^0(X)$ identifies
it with $\Z\beta_d$. To compute its degree, choose a point of $X/C_d$
away from the exceptional orbits. Its inverse image is one free
$C_d$-orbit and hence consists of exactly $d$ points. The map $q_d$ is
locally orientation-preserving at each of them, so it has topological
degree $d$, or equivalently
\[
 (q_d)_*[X]=d[X/C_d].
\]
Naturality of the first Chern class and of the pairing now
gives
\[
 \begin{aligned}
 \degc\beta_d
 &=\langle q_d^*c_1(\mathcal H),[X]\rangle\\
 &=\langle c_1(\mathcal H),(q_d)_*[X]\rangle=d.
 \end{aligned}
\]
Two classes with equal fiber data differ by an integer multiple of
$\beta_d$; equality of degrees makes that multiple zero.
\end{proof}

\subsection{Two convenient bases}

We now construct equivariant line bundles realizing the generators below. For each exceptional orbit $F$, perform the following construction and denote the resulting line bundle by $\lambda_F$. Choose one small disk around every point of $F$, with the disks pairwise disjoint and permuted by the group. Thus, when $F$ contains more than one point, the construction is carried out simultaneously on all these disks and produces a single line bundle $\lambda_F$. In a centered coordinate $w$, the stabilizer generator acts by $w\mapsto\zeta w$, with $\zeta=i$ at $\mathsf p,\mathsf q$ for $C_4$, and $\zeta=-1$ at every order-two center. On the trivial complex line bundle over each disk, define the stabilizer action by
\[
 (w,v)\longmapsto(\zeta w,\zeta v),
\]
and extend it to the other disks in the orbit by equivariance. Glue these bundles to the trivial equivariant line bundle on the complement of the centers as follows. Over an annular overlap, use $v_{-}$ for the fiber coordinate in the trivialization coming from the disk and $v_{+}$ for the fiber coordinate in the exterior trivialization. Since the center does not belong to the overlap, $w\ne0$ there, and $w^{-1}=1/w$ is a well-defined map from the annulus to $\C^\times=GL_1(\C)$. The transition function is
\begin{equation*}
 v_{+}=w^{-1}v_{-}.
\end{equation*}
Thus $w^{-1}$ is the reciprocal of the local complex coordinate. The transition is equivariant because the exterior fiber coordinate is fixed and
\[
 (\zeta w)^{-1}(\zeta v_{-})
 =w^{-1}v_{\mathrm{in}}=v_{+}.
\]
Finally, the constant exterior section $v_{+}=1$ is written in the disk trivialization as
\[
 1=w^{-1}v_{-},\qquad\text{hence}\qquad
 v_{-}=w.
\]
It therefore extends across the center and vanishes there. In the
oriented real coordinates underlying the complex coordinate $w$, its
local representative is the map $w\mapsto w$. Its derivative at the
center is the identity, so the zero is transverse, simple, and has local
index $+1$. Recall that, for a complex line bundle $L$ over an oriented
closed surface and a section with isolated zeros,
\[
 \langle c_1(L),[X]\rangle
 =\langle e(L_{\R}),[X]\rangle
 =\sum_{x\,:\,s(x)=0}\operatorname{ind}_x(s),
\]
where $L_{\R}$ is the underlying oriented real plane bundle; see
\cite[Theorem~11.17]{BT}. Repeating the construction equivariantly over
every disk in an orbit $F$ gives one positive simple zero at each point
of $F$ and no zeros on the exterior trivialization. Consequently the
resulting line bundle has degree $|F|$. At a center its stabilizer acts
on the fiber by the prescribed character $\zeta$, and equivariance
transports this character to every point of the orbit. At the other
exceptional points the exterior fiber coordinate is fixed, so their
fiber characters are trivial.

For $C_4$ call these bundles $\lambda_{\mathsf p}$, $\lambda_{\mathsf q}$ and $\lambda_R$, corresponding to $\{\mathsf p\}$, $\{\mathsf q\}$ and $\{x_{10},x_{01}\}$; their degrees are $1,1,2$. In particular, $\lambda_R$ is the one bundle obtained by carrying out the construction on the two disks centered at $x_{10}$ and $x_{01}$ simultaneously. For $C_2$ write $\lambda_{ab}$ for the degree-one bundle constructed separately at the singleton orbit $\{x_{ab}\}$, with $a,b\in\{0,1\}$. Set
\begin{align*}
 \mathcal A&=K_{C_4}^0(X),&
 P_k&=[\lambda_{\mathsf p}^{\otimes k}]-1,&
 Q_k&=[\lambda_{\mathsf q}^{\otimes k}]-1\quad (k=1,2,3),\\
 \mathcal B&=K_{C_2}^0(X),&
 U&=[\lambda_R]-1,&
 E_{ab}&=[\lambda_{ab}]-1.
\end{align*}
Then the ordered integral bases are
\begin{equation}\label{eq:bases}
 \mathcal A:\ (1,P_1,P_2,P_3,Q_1,Q_2,Q_3,U,\beta_4),\qquad
 \mathcal B:\ (1,E_{00},E_{10},E_{01},E_{11},\beta_2).
\end{equation}
Indeed $P_k,Q_k$ give the differences between the character with
value $i^k$ and the trivial character at their respective centers;
$U$ and each $E_{ab}$ give the nontrivial character minus the
trivial character at their respective order-two orbits. Together
with the common rank-one class, these form bases of
$\mathcal L_4,\mathcal L_2$. Adding $\beta_d$ gives the claimed
bases by Lemma~\ref{lem:surface-basis}. Their ranks are nine and six.
All classes except $1$ have rank zero, and their degrees are
\[
 \degc P_k=\degc Q_k=k,\qquad
 \degc U=2,\qquad \degc E_{ab}=1,\qquad \degc\beta_d=d.
\]


\subsection{Computation of the restriction map between the planar groups}\label{sec:restriction}

Under \eqref{eq:descent}, the inclusion $p2\hookrightarrow p4$ induces restriction of bundles
\[
 \rho=\Res_{C_2}^{C_4}\colon\mathcal A\longrightarrow\mathcal B.
\]

\begin{proposition}\label{prop:restriction}
In the bases \eqref{eq:bases},
\begin{equation*}
\begin{aligned}
 \rho(1)&=1,\\
 \rho(P_1)&=E_{00},&\rho(P_2)&=\beta_2,
     &\rho(P_3)&=E_{00}+\beta_2,\\
 \rho(Q_1)&=E_{11},&\rho(Q_2)&=\beta_2,
     &\rho(Q_3)&=E_{11}+\beta_2,\\
 \rho(U)&=E_{10}+E_{01},&\rho(\beta_4)&=2\beta_2.
\end{aligned}
\end{equation*}
In particular, its image is primitive of rank five and
\[
 \coker\rho=\Z\omega,\qquad \omega=[E_{10}]=-[E_{01}].
\]
\end{proposition}

\begin{proof}
At $\mathsf p$, restriction sends the character $i^k$ to $(-1)^k$.
The classes $\rho(P_1),\rho(P_2),\rho(P_3)$ have degrees $1,2,3$,
respectively, and the same exceptional fiber representations as
$E_{00}$, $\beta_2$, and $E_{00}+\beta_2$. The uniqueness statement in
Lemma~\ref{lem:surface-basis} therefore gives
\[
 \rho(P_1)=E_{00},\qquad
 \rho(P_2)=\beta_2,\qquad
 \rho(P_3)=E_{00}+\beta_2.
\]
The argument at $\mathsf q$ is identical.


The line bundle $\lambda_R$ has degree two and restricts to the nontrivial fiber character at both $x_{10}$ and $x_{01}$. The K-theory element $U$ consequently restricts to $E_{10}+E_{01}$. Finally $\rho(\beta_4)$ has trivial fiber classes and degree four;
it equals $2\beta_2$. 

The displayed formulas give
\begin{equation}\label{eq:image}
 \operatorname{im}\rho=
 \langle1,E_{00},E_{11},E_{10}+E_{01},\beta_2\rangle.
\end{equation}
But we can replace $E_{01}$ by $E_{10}+E_{01}$ in this basis, which proves the assertions about image and cokernel.
\end{proof}

Put $\mathcal N=\ker\rho$ and denote its inclusion into $\mathcal A$ by $\jmath$. A direct verification gives that $\mathcal N$ has the basis
\begin{equation}\label{eq:kernel-basis}
\begin{aligned}
 n_1&=P_3-P_1-P_2,& n_2&=Q_3-Q_1-Q_2,\\
 n_3&=P_2-Q_2,& n_4&=\beta_4-P_2-Q_2.
\end{aligned}
\end{equation}
%Indeed, a kernel vector has zero coefficients on $1,U$, opposite
%coefficients on $P_1,P_3$ and on $Q_1,Q_3$, and satisfies
%\[
 %\operatorname{coeff}_{P_2}+\operatorname{coeff}_{P_3}
 %+\operatorname{coeff}_{Q_2}+\operatorname{coeff}_{Q_3}
 %+2\operatorname{coeff}_{\beta_4}=0.
%\]
%Solving this last equation gives a unique integer combination
%of the four vectors in \eqref{eq:kernel-basis}.

\section{Equivariant \texorpdfstring{$K$}{K}-theory of
\texorpdfstring{$\underline{E}H$}{EH}}

\subsection{Actions on the bundle bases}

Let $a$ be the automorphism of $p4$ that fixes $u,v$ and sends $s$ to $us$. Translation by $\delta=(1/2,1/2)$ on the plane realizes this automorphism:
\[
 \tau_\delta h\tau_\delta^{-1}=a(h),\qquad
 (I-R)\delta=(1,0).
\]
To see this on the generators, write $u(x)=x+(1,0)$,
$v(x)=x+(0,1)$, $s(x)=Rx$, and
$\tau_\delta(x)=x+\delta$. Conjugation by a translation fixes every
translation, since translations commute; in particular,
$\tau_\delta u\tau_\delta^{-1}=u$ and
$\tau_\delta v\tau_\delta^{-1}=v$. For the rotation, direct evaluation
at $x\in\R^2$ gives
\[
 (\tau_\delta s\tau_\delta^{-1})(x)
 =R(x-\delta)+\delta
 =Rx+(I-R)\delta
 =Rx+(1,0)
 =(us)(x).
\]
Thus conjugation by $\tau_\delta$ has exactly the stated effect on the
generators of $p4$.
On $X$ this translation is an involution. The identity
$\tau_\delta(hx)=a(h)\tau_\delta(x)$ says that
$\tau_\delta$ is equivariant after the action on the target is
relabelled by $a$. More explicitly, if $\mathcal E\to\R^2$ is a
$p4$-equivariant vector bundle, its pullback has fiber
$(\tau_\delta^*\mathcal E)_x=\mathcal E_{\tau_\delta(x)}$. Give this
pullback the $p4$-action
\[
 h\cdot(x,e)=(hx,\,a(h)e).
\]
This is well-defined because $a(h)$ sends the fiber over
$\tau_\delta(x)$ to the fiber over $a(h)\tau_\delta(x)=\tau_\delta(hx)$.
Thus pulling back and relabelling the action defines a map on
equivariant vector bundles, and hence on their K-theory group
$K_{p4}^0(\R^2)\cong\mathcal A$; this is the automorphism $T$.
The same construction restricted to $p2$ gives $S$.
\begin{lemma}\label{lem:actions-on-bases}
On the bases \eqref{eq:bases}, the automorphisms $T$ and $S$ act by
\[
\begin{aligned}
 T\colon\mathcal A&\xrightarrow{\cong}\mathcal A,&
 P_k&\leftrightarrow Q_k,&
 1,U,\beta_4&\text{ fixed},\\
 S\colon\mathcal B&\xrightarrow{\cong}\mathcal B,&
 E_{00}&\leftrightarrow E_{11},\quad
 E_{10}\leftrightarrow E_{01},&
 1,\beta_2&\text{ fixed}.
\end{aligned}
\]
\end{lemma}
\begin{proof}
To read off these formulas, track the exceptional orbits, their fiber
characters, and the degrees. The map $\tau_\delta$ sends $\mathsf p$
to $\mathsf q$. The isotropy generator at $\mathsf p$ is $s$, while
the corresponding generator at $\mathsf q$ is $us=a(s)$. Thus a fiber
character $s\mapsto i^k$ is transported to the same character
$us\mapsto i^k$, not to its inverse. Since $\tau_\delta$ preserves
degree, this gives $T(P_k)=Q_k$; the reverse exchange similarly gives
$T(Q_k)=P_k$. The order-two orbit $\{x_{10},x_{01}\}$ is preserved as
a set, with its nontrivial fiber character unchanged, so the degree-two
class $U$ is fixed. For the $C_2$ action, the singleton orbits move as
$x_{00}\leftrightarrow x_{11}$ and $x_{10}\leftrightarrow x_{01}$.
Their nontrivial characters are preserved, giving the stated swaps of
the $E_{ab}$. The rank-one trivial class is fixed. Finally, each
$\beta_d$ has zero fiber classes and degree $d$; orientation
preservation keeps that degree unchanged, so $\beta_d$ is fixed. The
uniqueness statement in Lemma~\ref{lem:surface-basis} says that these
fiber data and degrees determine the classes and the proof is finished.
\end{proof}

Rotation by $s$ normalizes $C_2$ and therefore induces an
automorphism
\[
 Q\colon\mathcal B\xrightarrow{\cong}\mathcal B.
\]
It exchanges $E_{10},E_{01}$ and fixes the other four basis
elements. All three operators $Q$, $S$, and $T$ are involutions,
and
\begin{equation*}
 \rho T=S\rho,\qquad Q\rho=\rho,\qquad SQ=QS.
\end{equation*}
Since $T$ preserves $\mathcal N$, its restriction
$T_{\mathcal N}=T|_{\mathcal N}\colon\mathcal N\to\mathcal N$ is given
by
\begin{equation}\label{eq:kernel-action}
 n_1\leftrightarrow n_2,
 \qquad n_3\longmapsto-n_3,\qquad n_4\longmapsto n_4.
\end{equation}


\subsection{Mayer--Vietoris for the amalgam decomposition of
\texorpdfstring{$H$}{H}}

We order the vertex term so that the first copy of $\mathcal A$
corresponds to $V_0$ and the second copy corresponds to $V_1$.
Thus an element $(x,y)\in\mathcal A\oplus\mathcal A$ records the
restrictions to the two vertex spaces in that order.

Since the planar odd K-theory groups vanish, the Mayer--Vietoris sequence in equivariant $K$-theory associated with the decomposition $H=V_0*_E V_1$ gives
\begin{equation}\label{eq:MV-H}
 0\longrightarrow\KC{0}{H}
 \longrightarrow\mathcal A\oplus\mathcal A
 \xrightarrow{\partial}\mathcal B
 \longrightarrow\KC{1}{H}\longrightarrow0,
 \qquad \partial(x,y)=\rho x-\rho y.
\end{equation}
Equation~\eqref{eq:image} gives
\[
 \KC{1}{H}=\Z\omega.
\]
An explicit integral identification for the even group is
\begin{equation}\label{eq:H-coordinates}
 \mathcal A\oplus\mathcal N\xrightarrow{\cong}\ker\partial,
 \qquad (y,n)\longmapsto(y+\jmath n,y).
\end{equation}
Thus
\begin{equation}\label{eq:H-results}
 \KC{0}{H}\cong\Z^{13},\qquad \KC{1}{H}\cong\Z.
\end{equation}
\section{Computation of the \texorpdfstring{$K$}{K}-theory of
\texorpdfstring{$\Gamma_j$}{Gamma j}}
\subsection{The cohomological Wang sequence}

For $j=1,2$, define an affine map
$f_j\colon\underline EH\to\underline EH$ by
\[
 f_1(x,y)=(x+(1/2,1/2),y),\qquad
 f_2(x,y)=(x+(-1/2,1/2),y-1/2),
\]
where $x\in\R^2$ and $y\in\R$. These maps realize the automorphisms $\alpha_j$ in the sense that
\begin{equation*}
 f_j(hp)=\alpha_j(h)f_j(p),
 \qquad h\in H,\quad p\in\underline EH.
\end{equation*}
Let $t_j$ denote the generator of the cyclic quotient in
$\Gamma_j=H\rtimes_{\alpha_j}\langle t_j\rangle$. The action on
$\underline EH\times\R$ is defined on generators by
\[
 h\cdot(p,r)=(hp,r),\qquad
 t_j\cdot(p,r)=(f_j(p),r+1),
 \quad h\in H.
\]
\begin{comment}
In coordinates $p=(x,y)$ on $\underline EH=\R^2\times\R$, this says
\[
\begin{aligned}
 t_1(x,y,r)&=(x+(1/2,1/2),y,r+1),\\
 t_2(x,y,r)&=(x+(-1/2,1/2),y-1/2,r+1).
\end{aligned}
\]

We verify that these formulas respect the defining semidirect product
relation $t_jht_j^{-1}=\alpha_j(h)$. Identify $\underline EH$ with
$\R^3$ using coordinates $q=(x,y)$, and write $\tau_d(q)=q+d$ for
translation by $d$. Then $f_j=\tau_{d_j}$, where
\[
 d_1=(1/2,1/2,0),\qquad d_2=(-1/2,1/2,-1/2).
\]
The action of $s$ on $\underline EH$ is the linear map $q\mapsto Dq$,
where $D=\diag(R,-1)$. Since $t_j$ also translates the last factor
$\R$ by one, that translation cancels in a conjugation, and hence
\[
 (t_j h t_j^{-1})(q,r)
 =\bigl(f_j h f_j^{-1}(q),r\bigr).
\]
For $h=s$, the affine map on the right is
\[
 (f_jsf_j^{-1})(q)
 =d_j+D(q-d_j)
 =Dq+(I-D)d_j.
\]
The two translation vectors are
\[
 (I-D)d_1=(1,0,0),\qquad
 (I-D)d_2=(0,1,-1).
\]
These are the translation parts of $u$ and $vz^{-1}$, respectively.
Therefore
\[
 t_1st_1^{-1}=us,\qquad t_2st_2^{-1}=vz^{-1}s.
\]
For each translation generator $u,v,z$, conjugation by $f_j$ fixes it,
because translations commute with one another. Thus conjugation by $t_j$
agrees with $\alpha_j$ on the generators $u,v,z,s$ of $H$, and so
$t_jht_j^{-1}=\alpha_j(h)$ for every $h\in H$. This proves that the
formulas define a $\Gamma_j$-action. In general,
\[
 (h t_j^n)\cdot(p,r)=(h f_j^n(p),r+n),
 \qquad h\in H,\ n\in\Z.
\]
\end{comment}
The action is proper and cocompact and this is a model for $\underline E\Gamma_j$. Finally, the
map
\[
 \pi\colon\underline EH\times\R\longrightarrow S^1,\qquad
 \pi(p,r)=r\bmod\Z,
\]
is $\Gamma_j$-invariant: $H$ fixes $r$ and $t_j$ adds one. Cover $S^1$ by two open arcs
$I_0,I_1$ whose intersection is the disjoint union of two open arcs
$J_+,J_-$. Let $\mathcal U,\mathcal V$ be their inverse images in
$\underline EH\times\R$.

Choose a lift $\widetilde I_k\subset\R$ of each arc $I_k$, and put
$W_k=\underline EH\times\widetilde I_k$ for $k=0,1$. The preimage of
$I_k$ is the union of the integer translates of $W_k$. The subgroup
$H$ preserves $W_k$, and the translates by $t_j$ permute these copies
transitively. Consequently, as $\Gamma_j$-spaces,
\[
 \mathcal U\cong\Gamma_j\times_H W_0,\qquad
 \mathcal V\cong\Gamma_j\times_H W_1.
\]
Since each $\widetilde I_k$ is contractible and $H$ acts trivially on
that interval coordinate, $W_k$ is $H$-equivariantly homotopy equivalent
to $\underline EH$. Induction and homotopy invariance therefore give
\[
 K_{\Gamma_j}^i(\mathcal U)\cong\KC{i}{H},\qquad
 K_{\Gamma_j}^i(\mathcal V)\cong\KC{i}{H}.
\]
The intersection has two components, corresponding to $J_+$ and
$J_-$, and therefore
\[
 K_{\Gamma_j}^i(\mathcal U\cap\mathcal V)
 \cong\KC{i}{H}\oplus\KC{i}{H}.
\]
Choose the lifts $\widetilde I_0,\widetilde I_1$ so that on one
component of $I_0\cap I_1$, say $J_+$, the corresponding intervals in
$\R$ overlap in the same translate. The two local fiber coordinates
then use the same point of $\underline EH$, so the transition map on
that component is the identity. Choose them so that on the other
component $J_-$ the lift used for $I_1$ is the lift used for $I_0$
shifted by $+1$. To pass from the $I_0$ coordinate to the $I_1$
coordinate there, apply $t_j$; since
$t_j(p,r)=(f_j(p),r+1)$, the fiber coordinate changes by $f_j$. The induced map
on equivariant $K$-theory is $\alpha_j^*$, because
$f_j(hp)=\alpha_j(h)f_j(p)$. Thus, after identifying the two
intersection components with $\underline EH$ and choosing the usual
signs in the Mayer--Vietoris difference map, a pair $(x,y)$ restricts
to $x-y$ on $J_+$ and to $x-\alpha_j^*y$ on $J_-$. Therefore the
difference map is
\begin{equation*}
 D_i\colon\KC{i}{H}\oplus\KC{i}{H}
 \longrightarrow\KC{i}{H}\oplus\KC{i}{H},
 \qquad
 D_i(x,y)=(x-y,x-\alpha_j^*y).
\end{equation*}
Use $(x-y,y)$ as coordinates in the source and subtract the first
component of the target from the second. These integral changes of coordinates transform $D_i$ into
\[
 (u,y)\longmapsto\bigl(u,(1-\alpha_j^*)y\bigr).
\]
The identity summand contributes no kernel or cokernel. The Mayer--Vietoris sequence therefore reduces to the cohomological Wang sequence
\[
 \cdots\longrightarrow\KC{i}{\Gamma_j}
 \longrightarrow\KC{i}{H}
 \xrightarrow{1-\alpha_j^*}\KC{i}{H}
 \longrightarrow\KC{i+1}{\Gamma_j}\longrightarrow\cdots.
\]

Write
\[
 M_j=\alpha_j^*\big|_{\KC{0}{H}}.
\]

Then we must determine the action of $\alpha_j^*$ on $\KC{1}{H}$ and the kernels and cokernels of $1-M_j$. The next subsection computes the two induced actions, and Section~\ref{section:coinvariants} computes the required kernels and cokernels.

\subsection{The action of \texorpdfstring{$\alpha_j^*$}{alpha j}}

Extend the automorphism $a$ to $H$ by fixing $z$, and let $b$ be the automorphism of $H$ that fixes $u,v,z$ and sends $s$ to $zs$. Then
\[
 \alpha_1=a,\qquad \alpha_2=c_{u^{-1}}ab^{-1},
\]
where $c_{u^{-1}}$ denotes inner conjugation. Define the affine map
\[
 f_b(x,y)=(x,y+1/2).
\]
Conjugation by $f_1,f_b,f_2$ on the affine $H$-action gives $a,b,\alpha_2$, respectively. Inner conjugations induce the identity after the natural identification of equivariant bundles, so $c_{u^{-1}}$ does not change the induced maps on the groups in \eqref{eq:H-results}.

To describe the induced maps on the vertex term, put
\[
 \Pi_k=\R^2\times\{k/2\},\qquad
 V_k=\langle u,v,z^ks\rangle.
\]
Since $z^ks(x,y)=(Rx,k-y)$, the plane $\Pi_k$ has stabilizer $V_k$. In particular, $\Pi_0$ and $\Pi_1$ represent the two vertex orbits, with stabilizers $V_0$ and $V_1$. The map $f_1$ preserves both
planes and restricts on each of them to translation by
$(1/2,1/2)$ in the planar coordinate. It therefore induces $T$ on
each copy of $\mathcal A$.

On the other hand, $f_b(\Pi_k)=\Pi_{k+1}$ and $b(V_k)=V_{k+1}$. Thus $b$ exchanges the two vertex orbits.  Consequently, $b$ acts on the vertex term by interchanging the two copies of $\mathcal A$.

The direct sum $\mathcal A\oplus\mathcal A$ is the vertex term in the Mayer--Vietoris sequence \eqref{eq:MV-H}. Since $\alpha_1=a$ and $\alpha_2$ differs from $ab^{-1}$ by an inner automorphism, the latter first exchanges the two vertex classes and then applies $T$ to each. Define the resulting automorphisms by
\begin{equation}\label{eq:vertex-maps}
 \Theta_1,\Theta_2\colon\mathcal A\oplus\mathcal A
 \xrightarrow{\cong}\mathcal A\oplus\mathcal A,
 \qquad
 \Theta_1(x,y)=(Tx,Ty),\quad
 \Theta_2(x,y)=(Ty,Tx).
\end{equation}
Thus $\Theta_j$ records the action of $\alpha_j$ on the two vertex restrictions before passing to $K_H^0(\underline EH)$. The compatibility with the edge map $\partial$, verified below, shows that each $\Theta_j$ preserves $\ker\partial$. Under the identification $K_H^0(\underline EH)\cong\ker\partial$ from \eqref{eq:MV-H}, the restriction $\Theta_j|_{\ker\partial}$ is precisely $\alpha_j^*$ in even degree.

Let us determine the action of $\alpha_j^*$ for the odd group. Translating $[0,1/2]$ by $b$ gives $[1/2,1]$. The element $zs$ folds this interval back by $y\mapsto1-y$ and acts on the plane by $R$. Since $z$ acts trivially on the planar coordinate and rotation by $s$ induces $Q$ on $\mathcal B$, the planar part of this identification induces $Q$. The reversal of the edge orientation contributes an additional minus sign. Hence the induced endomorphism of the oriented edge term $\mathcal B$ is $-Q$. On the other hand, the planar translation defining $a$ induces $S$ on $\mathcal B$ by Lemma~\ref{lem:actions-on-bases}, and it preserves the orientation of the edge. Thus the induced endomorphism associated with $a$ is $S$. Recall that $\omega=[E_{10}]=-[E_{01}]$ is the generator of $\coker\rho\cong\KC{1}{H}$. Since
\[
 S\omega=-\omega,\qquad Q\omega=-\omega,
\]
the map $b$ acts as $+1$ on $\KC{1}{H}$, whereas $a$ acts as $-1$. It follows that
\begin{equation}\label{eq:odd-sign}
 \alpha_j^*\big|_{\KC{1}{H}}=-1,\qquad j=1,2.
\end{equation}
The compatibility of the vertex maps with the Mayer--Vietoris differential can be checked directly:
\[
 \partial\Theta_1=S\partial,\qquad
 \partial\Theta_2=(-QS)\partial.
\]
The second equality uses $Q\rho=\rho$. It induces the same sign $-1$ on the cokernel generated by $\omega$.

Under the identification $\KC{0}{H}\cong\ker\partial$, the map $M_j$ is the restriction of $\Theta_j$ to $\ker\partial$. Transporting this restriction through the isomorphism \eqref{eq:H-coordinates}, equation~\eqref{eq:vertex-maps} becomes
\begin{equation}\label{eq:monodromy-blocks}
 M_1=\begin{pmatrix}T&0\\0&T_{\mathcal N}\end{pmatrix},\qquad
 M_2=\begin{pmatrix}T&T\jmath\\0&-T_{\mathcal N}\end{pmatrix}.
\end{equation}
%For instance, $\Theta_2(y+\jmath n,y)=(Ty,Ty+T\jmath n)$,
%whose new coordinates are $(Ty+T\jmath n,-T_{\mathcal N}n)$.
%Both $M_j$ are integral involutions on a free group of rank thirteen.

\section{Coinvariants and the four-dimensional groups}\label{section:coinvariants}

The Wang sequence reduces the remaining calculation to the kernels and cokernels of $1-M_j$. We compute the cokernels first and then recover the kernel ranks.

For an automorphism $L$ of an abelian group $V$, write $V_L=V/(1-L)V$. Since $T$ interchanges three pairs in
\eqref{eq:bases},
\[
 \mathcal A_T\cong\Z^6,
 \qquad\text{with basis }1,P_1,P_2,P_3,U,\beta_4,
\]
where $Q_k=P_k$ in this quotient. Equation~\eqref{eq:kernel-action} gives
\[
 \mathcal N_{T_{\mathcal N}}\cong\Z^2\oplus\Z/2:
 \qquad n_1=n_2,\quad 2n_3=0,
\]
with $n_4$ unrestricted. Therefore
\begin{equation}\label{eq:coinv-first}
 \coker(1-M_1)\cong\Z^8\oplus\Z/2.
\end{equation}
%The torsion is represented by $n_3=P_2-Q_2$ in the kernel summand.

On the other hand, we will divide out the relations definig the cokernel in two steps. First to the first nine columns of $1-M_2$ we impose $(1-T)\mathcal A=0$. After quotienting by $\textcolor{black}{Img(1-T)}$ them, the remaining relations necessary to define $\coker(1-M_2)$, %\textcolor{red}{are given by the fact that for any $\overline{n}\in \mathcal{N}$
%\begin{equation}
%T(\overline{n})\oplus 0 = 0\oplus (1+T)(\overline{n}), \qquad\text{in coker}(1-M_{2}) 
%\end{equation}
%obtaining the relations
%\begin{equation*}
 %   [n_{1}\oplus 0] = [n_{2}\oplus 0]= [0\oplus n_{1}+n_{2}], \qquad [n_{3}\oplus 0]=0, \qquad [n_{4}\oplus 0]=[0\oplus 2n_{4}].
%\end{equation*}
%Therefore $[\mathcal{N}\oplus 0]\subseteq [0\oplus \mathcal{N}]\subseteq \coker(1-M_{2})$, which means that $\coker(1-M_{2})\cong \mathbb{Z}_{8}$ with base given by the set
%\begin{equation*}
%    [\{1,U,P_{1},P_{2}\}\oplus 0]\cup [0\oplus \{n_{1},n_{2},n_{3},n_{4}\}].
%\end{equation*} 
%}
indexed by $n\in\mathcal N$, are
\begin{equation}\label{eq:second-relations}
 \bigl(-[\jmath n],(1+T_{\mathcal N})n\bigr)=0
 \quad\text{in }\mathcal A_T\oplus\mathcal N.
\end{equation}
In $\mathcal A_T$ put
\[
 h=[P_3-P_1-P_2],\qquad k=[\beta_4-2P_2].
\]
The four kernel generators defined in \eqref{eq:kernel-basis} satisfy
\[
 [\jmath n_1]=[\jmath n_2]=h,\qquad
 [\jmath n_3]=0,\qquad [\jmath n_4]=k.
\]
Together with \eqref{eq:kernel-action}, this reduces \eqref{eq:second-relations} to exactly two relations in $\coker(1-M_2)$:
\begin{equation}\label{eq:two-relations}
 (-h,n_1+n_2)=0,\qquad (-k,2n_4)=0.
\end{equation}
The six elements
\[
 [1],[U],[P_1],[P_2],h,k
\]
are an integral basis of $\mathcal A_T$: replacing $P_3$ by $h$ and $\beta_4$ by $k$ is unimodular. Thus in $\coker(1-M_2)$ the relations \eqref{eq:two-relations} eliminate $h$ and $k$, leaving the eight free generators
\[
 [1],[U],[P_1],[P_2],[n_1],[n_2],[n_3],[n_4].
\]
In particular,
\begin{equation*}
 \coker(1-M_2)\cong\Z^8.
\end{equation*}

Both $1-M_j$ are endomorphisms of $\Z^{13}$. The ranks of
their kernels and cokernels therefore agree, so
\begin{equation}\label{eq:invariant-ranks}
 \ker(1-M_j)\cong\Z^8,\qquad j=1,2.
\end{equation}

By \eqref{eq:odd-sign}, the odd-degree middle map is
multiplication by two on $\Z$. Its kernel is zero. Exactness gives
\begin{align}
 0\longrightarrow\Z/2\longrightarrow\KC{0}{\Gamma_j}
 &\longrightarrow\ker(1-M_j)\longrightarrow0,
 \label{eq:wang-even}\\
 \KC{1}{\Gamma_j}&\cong\coker(1-M_j).
 \label{eq:wang-odd}
\end{align}
The right-hand term of \eqref{eq:wang-even} is free by \eqref{eq:invariant-ranks}; hence this sequence splits as a sequence of abelian groups. Combining \eqref{eq:coinv-first}--\eqref{eq:wang-odd} proves
\eqref{eq:cohom-results}.

\subsection{Reduced group \texorpdfstring{$C^*$}{C-star}-algebra
\texorpdfstring{$K$}{K}-theory}

%\subsection{From proper $K$-cohomology to $K$-homology}

For a finite proper $G$-CW-complex $Y$, the equivariant universal coefficient theorem of Joachim--L\"uck \cite[Theorem~0.3, equation~(0.5)]{JL} gives
\begin{equation}\label{eq:UCT}
 0\longrightarrow
 \operatorname{Ext}^1_{\Z}(K_G^{i+1}(Y),\Z)
 \longrightarrow K_i^G(Y)
 \longrightarrow\operatorname{Hom}_{\Z}(K_G^i(Y),\Z)
 \longrightarrow0.
\end{equation}
The hypotheses hold for $Y=\underline E\Gamma_j=\R^4$.

For a finitely generated group $A\cong\Z^r\oplus T$ with $T$
finite,
\[
 \operatorname{Hom}_{\Z}(A,\Z)\cong\Z^r,
 \qquad\operatorname{Ext}^1_{\Z}(A,\Z)
 \cong\operatorname{Hom}(T,\Q/\Z)\cong T.
\]
Since the right-hand term in \eqref{eq:UCT} is free, the sequence splits as a sequence of abelian groups. Thus the free ranks retain their parity, while the torsion comes from the opposite cohomological parity. Applying \eqref{eq:UCT} to \eqref{eq:cohom-results} proves \eqref{eq:hom-results}.

Average an inner product over the finite holonomy group. The original affine action of each $\Gamma_j$ on $\R^4$ is then proper and isometric. The Higson--Kasparov theorem \cite{HK} gives the Baum--Connes assembly isomorphism
\[
 \mu_{\Gamma_j}\colon K_i^{\Gamma_j}(\R^4)
 \xrightarrow{\cong}K_i(C_r^*(\Gamma_j)),\qquad i=0,1.
\]
This transfers the second table in Theorem~\ref{thm:main} to the reduced group algebras. 

\begin{remark}
The same groups could alternatively be obtained by computing Bredon homology with coefficients in complex representation-ring and then using the equivariant Atiyah--Hirzebruch spectral sequence for equivariant $K$-homology. That approach requires keeping track of the resulting filtration and its additive extensions. We have preferred the more direct $K$-theoretic method above, since the planar restriction map and the two monodromy actions can be described explicitly.
\end{remark}

\begin{thebibliography}{99}
\bibitem{BT}
R. Bott and L. W. Tu,
\emph{Differential forms in algebraic topology},
Graduate Texts in Mathematics, vol.~82, Springer-Verlag,
New York--Berlin, 1982.

\bibitem{DL}
J. F. Davis and W. L\"uck,
\emph{The topological $K$-theory of certain crystallographic groups},
J. Noncommut. Geom. \textbf{7} (2013), no.~2,
373--431.

\bibitem{HK}
N. Higson and G. Kasparov,
\emph{$E$-theory and $KK$-theory for groups which act properly and
isometrically on Hilbert space},
Invent. Math. \textbf{144} (2001), 23--74.

\bibitem{JL}
M. Joachim and W. L\"uck,
\emph{Topological $K$-(co-)homology of classifying spaces of discrete groups},
Algebr. Geom. Topol. \textbf{13} (2013), no.~1, 1--34.

\bibitem{LL}
M. Langer and W. L\"uck,
\emph{Topological $K$-theory of the group $C^*$-algebra of a
semi-direct product $\Z^n\rtimes\Z/m$ for a free conjugation action},
J. Topol. Anal. \textbf{4} (2012), no.~2, 121--172.

\bibitem{LO}
W. L\"uck and B. Oliver,
\emph{The completion theorem in $K$-theory for proper actions of a discrete group},
Topology \textbf{40} (2001), 585--616.

\bibitem{LuckBook}
W. L\"uck,
\emph{Isomorphism conjectures in $K$- and $L$-theory},
Ergebnisse der Mathematik und ihrer Grenzgebiete, 3.~Folge,
vol.~78, Springer, Cham, 2025.

\bibitem{LuckSurvey}
W. L\"uck,
\emph{Survey on classifying spaces for families of subgroups},
in \emph{Infinite groups: geometric, combinatorial and dynamical
aspects}, Progress in Mathematics, vol.~248, Birkh\"auser, Basel,
2005, 269--322.

\bibitem{PP}
N. Petrosyan and B. Putrycz,
\emph{On cohomology of crystallographic groups with cyclic holonomy of split type},
J. Algebra \textbf{367} (2012), 237--246.

\bibitem{SV}
L. J. S\'anchez Salda\~na and M. Vel\'asquez,
\emph{On the $K$-theory of groups of the form
$\Z^n\rtimes\Z/m$ with $m$ square-free},
Proc. Amer. Math. Soc. \textbf{155} (2027), no. 1, 127--140.
\end{thebibliography}
\end{document}
